\documentclass[11pt]{article}

\usepackage[margin=1in]{geometry}
\usepackage[T1]{fontenc}
\usepackage[utf8]{inputenc}
\usepackage{lmodern}
\usepackage{amsmath,amssymb,mathtools}
\usepackage{booktabs}
\usepackage{enumitem}
\usepackage[hidelinks]{hyperref}
\usepackage{microtype}

\title{SimpleSFT Structural Memory Estimator Report}
\author{}
\date{\today}

\newcommand{\bytes}[1]{\operatorname{bytes}\!\left(#1\right)}
\newcommand{\ceil}[1]{\left\lceil #1 \right\rceil}
\newcommand{\floor}[1]{\left\lfloor #1 \right\rfloor}

\begin{document}
\maketitle

\section{Purpose}
This document describes the \emph{current implemented estimator} in
SimpleSFT. It is not a report about an older calibration-heavy model, nor a
wish list of future refinements. Every formula here is intended to match the
current code paths in:
\begin{itemize}[leftmargin=1.5em]
\item \texttt{resident\_state\_model.py},
\item \texttt{activation\_model.py},
\item \texttt{workspace\_model.py},
\item \texttt{phase\_model.py},
\item \texttt{parallelism\_model.py},
\item \texttt{optimizer\_model.py},
\item \texttt{estimate.py}.
\end{itemize}

The estimator is designed to be:
\begin{itemize}[leftmargin=1.5em]
\item structural rather than calibration-driven,
\item explicit about saved tensors and workspace windows,
\item compositional across tuning mode, ZeRO mode, tensor parallelism, and
sequence parallelism,
\item simple enough that each large term can be traced to a named object.
\end{itemize}

The primary public estimate is an \emph{allocated per-rank peak}. The result
also carries a fixed backend runtime floor and rich debug metadata, but the
headline estimate is not an allocator-cache fit.

\section{High-Level Result}
For a given model and estimator config, the estimator constructs:
\begin{enumerate}[leftmargin=1.5em]
\item resident state,
\item retained activations,
\item workspace and communication windows,
\item synthetic phase peaks.
\end{enumerate}

The three main phase peaks are:
\begin{align}
M_{\mathrm{forward}} &= M_{\mathrm{resident,base}}
+ M_{\mathrm{act,fwd}}
+ M_{\mathrm{ws,fwd}}, \\
M_{\mathrm{backward}} &= M_{\mathrm{resident,with\ grad}}
+ M_{\mathrm{act,bwd}}
+ M_{\mathrm{ws,bwd}}, \\
M_{\mathrm{optimizer}} &= M_{\mathrm{resident,with\ grad}}
+ M_{\mathrm{bwd,end}}
+ M_{\mathrm{ws,opt}}.
\end{align}

The global peak is then:
\begin{equation}
M_{\mathrm{peak}} =
\max\!\left(
M_{\mathrm{forward}},
M_{\mathrm{backward}},
M_{\mathrm{optimizer}}
\right).
\end{equation}

Feasibility is:
\begin{equation}
M_{\mathrm{peak}} \leq M_{\mathrm{gpu}},
\end{equation}
where $M_{\mathrm{gpu}}$ is the per-rank GPU capacity in bytes.

\section{Inputs and Structural Objects}
\subsection{EstimatorConfig}
The estimator accepts \texttt{EstimatorConfig}. The relevant structural fields
are:
\begin{itemize}[leftmargin=1.5em]
\item tuning mode: \texttt{full\_ft} or \texttt{lora},
\item distributed mode: \texttt{single\_gpu}, \texttt{ddp}, \texttt{zero2},
\texttt{zero3},
\item optimizer family,
\item attention backend,
\item checkpointing flag,
\item tensor parallel degree,
\item sequence parallel flag,
\item vocab-parallel logits flag,
\item sequence length and micro-batch size,
\item topology: \texttt{gpus\_per\_node}, \texttt{num\_nodes},
\item dtype choices,
\item LoRA config,
\item optional runtime-floor override.
\end{itemize}

\subsection{ModelSpec}
The estimator uses a \texttt{ModelSpec}. The important fields are:
\begin{itemize}[leftmargin=1.5em]
\item dense model shape: number of layers, hidden size, intermediate size,
attention heads, vocab size,
\item total parameter count,
\item trainable linear layer summaries,
\item optional per-parameter summaries.
\end{itemize}

If parameter summaries are present, resident-state accounting is done
parameter-by-parameter. If they are absent, some paths fall back to coarse
total-parameter accounting.

\subsection{Notation}
We use:
\begin{align}
L &= \text{number of transformer layers}, &
H &= \text{hidden size}, \\
I &= \text{intermediate size}, &
V &= \text{vocabulary size}, \\
B &= \text{micro-batch size per rank}, &
S &= \text{sequence length}, \\
T &= B \cdot S, &
W &= \text{world size}, \\
T_{\mathrm{TP}} &= \text{tensor parallel degree}, &
D_{\mathrm{P}} &= \text{data parallel degree after TP}.
\end{align}

Byte widths:
\begin{align}
b_w &= \bytes{\text{weight dtype}}, &
b_g &= \bytes{\text{gradient dtype}}, \\
b_o &= \bytes{\text{optimizer-state dtype}}, &
b_u &= \bytes{\text{optimizer-update dtype}}, \\
b_m &= \bytes{\text{master-weight dtype}}, &
b_a &= \bytes{\text{adapter-weight dtype}}, \\
b_{ag} &= \bytes{\text{adapter-gradient dtype}}, &
b_{ao} &= \bytes{\text{adapter optimizer-state dtype}}.
\end{align}

\section{Parallelism Semantics}
\subsection{World Size, TP, and DP}
The estimator defines:
\begin{align}
W &=
\begin{cases}
1 & \text{for single\_gpu}, \\
\texttt{num\_nodes} \cdot \texttt{gpus\_per\_node} & \text{otherwise},
\end{cases} \\
T_{\mathrm{TP}} &= \texttt{tensor\_parallel\_degree}, \\
D_{\mathrm{P}} &=
\begin{cases}
1 & \text{for single\_gpu}, \\
\dfrac{W}{T_{\mathrm{TP}}} & \text{otherwise}.
\end{cases}
\end{align}

Validity requires:
\begin{itemize}[leftmargin=1.5em]
\item $T_{\mathrm{TP}} \geq 1$,
\item $W \bmod T_{\mathrm{TP}} = 0$,
\item single-GPU mode may not use TP,
\item sequence parallel requires $T_{\mathrm{TP}} > 1$.
\end{itemize}

\subsection{Sequence Parallel}
Sequence parallel is \emph{not} an independent degree in the current model. It
is a boolean mode that reuses the TP group. The sequence-local divisor is:
\begin{equation}
D_{\mathrm{SP}} =
\begin{cases}
1 & \text{if sequence parallel is off or TP is off}, \\
T_{\mathrm{TP}} & \text{if sequence parallel is on}.
\end{cases}
\end{equation}

Therefore the local token count used for sequence-local tensors is:
\begin{equation}
T_{\mathrm{local}} = \ceil{\dfrac{T}{D_{\mathrm{SP}}}}.
\end{equation}

\subsection{TP Sharding Heuristics}
The current estimator treats only certain tensors as TP-sharded:
\begin{itemize}[leftmargin=1.5em]
\item attention and MLP matrices,
\item embeddings and LM head when vocab-parallel logits are enabled.
\end{itemize}

This decision is made from parameter categories and a small row-parallel module
name heuristic, not from a full Megatron graph trace. That is accurate enough
for the estimator's structural goals, but it is still a heuristic.

\section{Trainable Parameter Set}
\subsection{Full Fine-Tuning}
For \texttt{full\_ft}, the trainable parameter set is:
\begin{equation}
\mathcal{P}_{\mathrm{train}} = \mathcal{P}_{\mathrm{all}},
\end{equation}
or, if explicit parameter specs are unavailable, one synthetic parameter tensor
of size $P$.

\subsection{LoRA}
For \texttt{lora}, the estimator synthesizes trainable adapter tensors from the
targeted linear layers. For each targeted linear layer with input width
$d_{\mathrm{in}}$ and output width $d_{\mathrm{out}}$, rank $r$, and optional
bias:
\begin{align}
\mathrm{LoRA\ A} &: (r, d_{\mathrm{in}}), \\
\mathrm{LoRA\ B} &: (d_{\mathrm{out}}, r), \\
\mathrm{bias} &: (d_{\mathrm{out}}) \quad \text{if enabled.}
\end{align}

These synthetic LoRA tensors are the trainable parameter set for LoRA tuning.

\section{Resident State}
\subsection{Local Parameter Bytes}
For a parameter tensor with total element count $n$, TP divisor
$d_{\mathrm{tp}}$, and optional ZeRO sharding:
\begin{equation}
n_{\mathrm{local,tp}} = \ceil{\dfrac{n}{d_{\mathrm{tp}}}}.
\end{equation}

If a resident tensor is ZeRO-sharded, the estimator further applies:
\begin{equation}
n_{\mathrm{local}} = \ceil{\dfrac{n_{\mathrm{local,tp}}}{D_{\mathrm{P}}}}.
\end{equation}

The generic local-byte rule is then:
\begin{equation}
M = n_{\mathrm{local}} \cdot \bytes{\mathrm{dtype}}.
\end{equation}

\subsection{Parameters}
If full parameter specs are available, the estimator sums local bytes over all
parameter tensors:
\begin{equation}
M_{\mathrm{params}} =
\sum_{p \in \mathcal{P}_{\mathrm{base}}} M_p
+ \mathbb{1}_{\mathrm{lora}}
\sum_{a \in \mathcal{P}_{\mathrm{adapter}}} M_a.
\end{equation}

The base-model parameter term uses:
\begin{itemize}[leftmargin=1.5em]
\item TP sharding when the parameter category is TP-shardable,
\item additional ZeRO sharding only for \texttt{zero3}.
\end{itemize}

For LoRA, adapter parameters are also sharded under \texttt{zero3}.

\paragraph{Fallback path.}
If explicit parameter specs are missing, the estimator falls back to:
\begin{equation}
M_{\mathrm{params}} \approx P \cdot b_w
\end{equation}
for the dense model. This fallback is intentionally coarse and cannot express
TP/ZeRO structure as accurately as the parameter-spec path.

\subsection{Gradients}
The resident gradient term is built from the trainable parameter set:
\begin{equation}
M_{\mathrm{grads}} =
\sum_{p \in \mathcal{P}_{\mathrm{train}}}
\ceil{
\dfrac{\ceil{n_p / d_{\mathrm{tp}}}}{d_{\mathrm{zero}}}
}
\cdot b_g^\star,
\end{equation}
where:
\begin{equation}
d_{\mathrm{zero}} =
\begin{cases}
D_{\mathrm{P}} & \text{for zero2 or zero3}, \\
1 & \text{otherwise},
\end{cases}
\end{equation}
and:
\begin{equation}
b_g^\star =
\begin{cases}
b_{ag} & \text{for LoRA}, \\
b_g & \text{for full FT}.
\end{cases}
\end{equation}

\subsection{Optimizer State}
For each trainable parameter $p$, the optimizer-state element count is:
\begin{equation}
n_{\mathrm{state}}(p) =
\begin{cases}
2 n_p & \text{for Adam or AdamW}, \\
0 & \text{for SGD}, \\
n_p & \text{for RMSprop}, \\
n_p & \text{for Adagrad}, \\
\mathrm{Adafactor}(p) & \text{for Adafactor}.
\end{cases}
\end{equation}

For Adafactor:
\begin{equation}
\mathrm{Adafactor}(p) =
\begin{cases}
n_p & \text{if } p \text{ is not matrix-like}, \\
\text{rows}(p) + \text{cols}(p) & \text{if } p \text{ is matrix-like}.
\end{cases}
\end{equation}

The local TP-adjusted state element count is formed by scaling the local
parameter count by the state ratio:
\begin{equation}
n_{\mathrm{state,local}}(p) =
\ceil{
\left(
\dfrac{n_{\mathrm{state}}(p)}{\max(n_p, 1)}
\right)
\cdot
\ceil{\dfrac{n_p}{d_{\mathrm{tp}}}}
}.
\end{equation}

State bytes are then:
\begin{equation}
M_{\mathrm{opt}} =
\begin{cases}
\left(\sum_p n_{\mathrm{state,local}}(p)\right) b_o^\star &
\text{for single\_gpu or ddp}, \\
\ceil{
\dfrac{
\left(\sum_p n_{\mathrm{state,local}}(p)\right) b_o^\star
}{D_{\mathrm{P}}}
} &
\text{for zero2 or zero3},
\end{cases}
\end{equation}
with resolved optimizer-state dtype:
\begin{equation}
b_o^\star =
\begin{cases}
\bytes{\texttt{optimizer\_state\_dtype}} &
\text{if explicitly set}, \\
\bytes{\texttt{fp32}} &
\text{if dtype is auto and ZeRO is on}, \\
b_{ao} & \text{if dtype is auto and LoRA is used}, \\
b_w & \text{otherwise}.
\end{cases}
\end{equation}

\subsection{Master Weights}
If master weights are enabled:
\begin{equation}
M_{\mathrm{master}} =
\sum_{p \in \mathcal{P}_{\mathrm{train}}}
\ceil{
\dfrac{\ceil{n_p / d_{\mathrm{tp}}}}{d_{\mathrm{zero}}}
}
\cdot b_m,
\end{equation}
where $d_{\mathrm{zero}} = D_{\mathrm{P}}$ for ZeRO and $1$ otherwise.

\subsection{Runtime Floor}
The runtime floor is a fixed backend constant unless overridden:
\begin{equation}
M_{\mathrm{runtime}} =
\begin{cases}
0.30~\mathrm{GiB} & \text{single\_gpu}, \\
0.60~\mathrm{GiB} & \text{ddp}, \\
0.80~\mathrm{GiB} & \text{zero2}, \\
1.00~\mathrm{GiB} & \text{zero3}.
\end{cases}
\end{equation}

\subsection{Resident Floors by Phase}
The phase model uses two resident floors:
\begin{align}
M_{\mathrm{resident,base}} &=
M_{\mathrm{params}} + M_{\mathrm{opt}} + M_{\mathrm{master}} + M_{\mathrm{runtime}},
\\
M_{\mathrm{resident,with\ grad}} &=
M_{\mathrm{resident,base}} + M_{\mathrm{grads}}.
\end{align}

\section{Retained Activations}
\subsection{Target Layer Set}
The activation model first selects a target set of trainable linear layers:
\begin{equation}
\mathcal{L}_{\mathrm{target}} =
\begin{cases}
\mathcal{L}_{\mathrm{trainable}} & \text{for full FT}, \\
\{ \ell \in \mathcal{L}_{\mathrm{trainable}} :
\ell.\mathrm{name} \text{ ends with a LoRA target suffix} \}
& \text{for LoRA}.
\end{cases}
\end{equation}

\paragraph{Implementation note.}
The helper named \texttt{\_average\_block\_layers()} currently just returns the
stored trainable linear layer tuple. The current model therefore treats the
\texttt{trainable\_linear\_layers} list as the representative per-layer/block
linear inventory.

\subsection{Checkpoint Segment Size}
The current implementation hardcodes:
\begin{equation}
L_{\mathrm{seg}} = 1.
\end{equation}

Then:
\begin{equation}
L_{\mathrm{act}} =
\begin{cases}
L & \text{if checkpointing is off}, \\
L_{\mathrm{seg}} & \text{if checkpointing is on}.
\end{cases}
\end{equation}

\subsection{Local Linear Widths Under TP}
For TP-sharded attention/MLP layers, the estimator uses local widths:
\begin{align}
d_{\mathrm{in,local}} &=
\begin{cases}
d_{\mathrm{in}} & \text{if not TP-sharded}, \\
d_{\mathrm{in}} & \text{if column-parallel}, \\
\ceil{d_{\mathrm{in}} / T_{\mathrm{TP}}} & \text{if row-parallel},
\end{cases}
\\
d_{\mathrm{out,local}} &=
\begin{cases}
d_{\mathrm{out}} & \text{if not TP-sharded}, \\
\ceil{d_{\mathrm{out}} / T_{\mathrm{TP}}} & \text{if column-parallel}, \\
d_{\mathrm{out}} & \text{if row-parallel}.
\end{cases}
\end{align}

This rule is used for saved linear inputs and hook-visible diagnostic outputs.

\subsection{Saved Activation Classes}
The activation model uses these explicit classes.

\paragraph{Saved linear inputs.}
\begin{equation}
M_{\mathrm{saved\_linear}} =
T_{\mathrm{local}}
\cdot L_{\mathrm{act}}
\cdot
\left(
\sum_{\ell \in \mathcal{L}_{\mathrm{target}}}
d_{\mathrm{in,local}}(\ell)
\right)
\cdot b_w.
\end{equation}

\paragraph{Residual and norm state.}
\begin{equation}
M_{\mathrm{residual}} =
T_{\mathrm{local}}
\cdot L_{\mathrm{act}}
\cdot 2H
\cdot b_w.
\end{equation}

\paragraph{Checkpoint boundaries.}
\begin{equation}
M_{\mathrm{ckpt}} =
T_{\mathrm{local}}
\cdot (L + 1)
\cdot H
\cdot b_w.
\end{equation}

\paragraph{Expanded-query saved context.}
Let the model's representative \texttt{q\_proj} and \texttt{o\_proj} widths be:
\begin{equation}
d_{q,\mathrm{out}}, \qquad d_{o,\mathrm{in}}.
\end{equation}
If $d_{q,\mathrm{out}} \leq H$, the term is zero. Otherwise:
\begin{equation}
M_{\mathrm{expand}} =
T_{\mathrm{local}}
\cdot L_{\mathrm{act}}
\cdot
\ceil{
\dfrac{
(d_{q,\mathrm{out}} - H)
+ \max(d_{o,\mathrm{in}} - H, 0)
}{T_{\mathrm{TP}}}
}
\cdot b_w.
\end{equation}

\paragraph{Attention-saved state.}
Only the \texttt{standard} backend retains an explicit quadratic score object:
\begin{equation}
M_{\mathrm{attn}} =
\begin{cases}
B \cdot
\ceil{\dfrac{\text{heads}}{T_{\mathrm{TP}}}}
\cdot S^2
\cdot L
\cdot \bytes{\mathrm{fp32}}
& \text{for standard attention}, \\
0 & \text{for sdpa or flash2}.
\end{cases}
\end{equation}

\paragraph{Loss state.}
\begin{equation}
M_{\mathrm{loss}} =
T_{\mathrm{local}}
\cdot
\ceil{
\dfrac{V}{d_{\mathrm{vocab}}}
}
\cdot b_w,
\end{equation}
where:
\begin{equation}
d_{\mathrm{vocab}} =
\begin{cases}
T_{\mathrm{TP}} & \text{if vocab-parallel logits are on}, \\
1 & \text{otherwise}.
\end{cases}
\end{equation}

\paragraph{LoRA low-rank intermediates.}
If LoRA is disabled, the term is zero. Otherwise:
\begin{equation}
M_{\mathrm{lora}} =
T_{\mathrm{local}}
\cdot L_{\mathrm{act}}
\cdot
\left(
\sum_{\ell \in \mathcal{L}_{\mathrm{target}}}
r_{\mathrm{local}}(\ell)
\right)
\cdot b_a,
\end{equation}
where:
\begin{equation}
r_{\mathrm{local}}(\ell) =
\begin{cases}
r & \text{if TP is off}, \\
\ceil{r / T_{\mathrm{TP}}} & \text{if TP is on}.
\end{cases}
\end{equation}

\paragraph{Hook-visible outputs.}
The estimator also computes:
\begin{equation}
M_{\mathrm{hook}} =
T_{\mathrm{local}}
\cdot
\left(
\sum_{\ell \in \mathcal{L}_{\mathrm{trainable}}}
d_{\mathrm{out,local}}(\ell)
\right)
\cdot b_w,
\end{equation}
but this is \emph{diagnostic only}. It is never added to retained activation
bytes.

\subsection{Forward and Backward Activation Formulas}
\paragraph{Checkpointing on.}
\begin{align}
M_{\mathrm{act,fwd}} &=
M_{\mathrm{ckpt}}
+ M_{\mathrm{attn}}
+ M_{\mathrm{loss}},
\\
M_{\mathrm{act,bwd}} &=
M_{\mathrm{ckpt}}
+ M_{\mathrm{saved\_linear}}
+ M_{\mathrm{residual}}
+ M_{\mathrm{attn}}
+ M_{\mathrm{lora}}
+ M_{\mathrm{expand}}.
\end{align}

\paragraph{Checkpointing off.}
\begin{align}
M_{\mathrm{act,bwd}} &=
M_{\mathrm{saved\_linear}}
+ M_{\mathrm{residual}}
+ M_{\mathrm{attn}}
+ M_{\mathrm{lora}}
+ M_{\mathrm{expand}},
\\
M_{\mathrm{act,fwd}} &=
M_{\mathrm{act,bwd}} + M_{\mathrm{loss}}.
\end{align}

\section{Workspace and Communication Windows}
\subsection{Forward Attention Workspace}
For \texttt{standard} attention, the estimator uses a tile-based forward
workspace:
\begin{equation}
M_{\mathrm{attn,fwd}} =
B \cdot
\ceil{\dfrac{\text{heads}}{T_{\mathrm{TP}}}}
\cdot S
\cdot \min(S, \mathrm{tile})
\cdot L
\cdot \bytes{\mathrm{fp32}},
\end{equation}
where:
\begin{equation}
\mathrm{tile} =
\begin{cases}
1024 & \text{standard}, \\
256 & \text{sdpa}, \\
128 & \text{flash2}.
\end{cases}
\end{equation}

For \texttt{sdpa} and \texttt{flash2}, the estimator instead uses a hidden-size
workspace:
\begin{equation}
M_{\mathrm{attn,fwd}} =
T_{\mathrm{local}}
\cdot \ceil{\dfrac{H}{T_{\mathrm{TP}}}}
\cdot L
\cdot b_w.
\end{equation}

\subsection{Backward Kernel Workspace}
The backward kernel workspace is:
\begin{equation}
M_{\mathrm{bwd,kernel}} =
M_{\mathrm{attn,fwd}}
+ T_{\mathrm{local}}
\cdot \ceil{\dfrac{I}{T_{\mathrm{TP}}}}
\cdot b_w.
\end{equation}

\subsection{Checkpoint Recompute Workspace}
If checkpointing is off:
\begin{equation}
M_{\mathrm{recompute}} = 0.
\end{equation}

If checkpointing is on:
\begin{equation}
M_{\mathrm{recompute}} =
T_{\mathrm{local}}
\cdot L_{\mathrm{seg}}
\cdot
\left(
H + \ceil{\dfrac{I}{T_{\mathrm{TP}}}}
\right)
\cdot b_w.
\end{equation}

\subsection{TP and SP Communication Windows}
If TP is active:
\begin{equation}
M_{\mathrm{tp,comm}} =
T_{\mathrm{local}} \cdot H \cdot b_w.
\end{equation}

If SP is also active:
\begin{equation}
M_{\mathrm{sp,comm}} =
T_{\mathrm{local}} \cdot H \cdot b_w.
\end{equation}

Otherwise each corresponding term is zero.

\subsection{Optimizer Update Scratch}
The per-parameter optimizer update element count is:
\begin{equation}
n_{\mathrm{update}}(p) = n_p
\end{equation}
for all currently supported optimizers:
\begin{equation}
\{\mathrm{adam},\ \mathrm{adamw},\ \mathrm{sgd},\ \mathrm{rmsprop},\ \mathrm{adagrad},\ \mathrm{adafactor}\}.
\end{equation}

So:
\begin{equation}
N_{\mathrm{update}} =
\sum_{p \in \mathcal{P}_{\mathrm{train}}} n_p.
\end{equation}

The resolved update dtype is:
\begin{equation}
b_u^\star =
\begin{cases}
\bytes{\texttt{optimizer\_update\_dtype}} &
\text{if explicitly set}, \\
b_{ao} &
\text{if dtype is auto, ZeRO is on, and LoRA is used}, \\
\bytes{\mathrm{fp32}} &
\text{if dtype is auto, ZeRO is on, and full FT is used}, \\
b_a &
\text{if dtype is auto, LoRA is used without ZeRO and optimizer is state-like}, \\
b_w &
\text{if dtype is auto, full FT is used without ZeRO and optimizer is state-like}, \\
b_{ag} &
\text{if dtype is auto, LoRA is used and optimizer is gradient-like}, \\
b_g &
\text{if dtype is auto, full FT is used and optimizer is gradient-like}.
\end{cases}
\end{equation}

The current implementation then forms:
\begin{equation}
M_{\mathrm{update,scratch}} =
\begin{cases}
N_{\mathrm{update}} \cdot b_u^\star &
\text{for single\_gpu or ddp}, \\
\ceil{\dfrac{N_{\mathrm{update}}}{W}} \cdot b_u^\star &
\text{for zero2 or zero3}.
\end{cases}
\end{equation}

\paragraph{Important note.}
This ZeRO scratch sharding is the \emph{current implemented behavior}. It is a
user-directed structural assumption and should not be confused with a
corpus-fitted claim about true runtime overlap.

\subsection{Bucket and Window Constants}
The estimator uses fixed bucket caps:
\begin{align}
M_{\mathrm{bucket,max}} &= 512~\mathrm{MiB}, \\
M_{\mathrm{prefetch,max}} &= 256~\mathrm{MiB}.
\end{align}

The bucket helper is:
\begin{equation}
\mathrm{bucket}(x) = \min(x, 512~\mathrm{MiB}).
\end{equation}

\subsection{DDP Windows}
For DDP:
\begin{align}
M_{\mathrm{ddp,bucket}} &= \mathrm{bucket}(M_{\mathrm{grads}}), \\
M_{\mathrm{ddp,comm}} &= M_{\mathrm{ddp,bucket}} + M_{\mathrm{grads}}.
\end{align}

\subsection{ZeRO Windows}
For \texttt{zero2} and \texttt{zero3}:
\begin{align}
M_{\mathrm{zero,fetch}} &=
M_{\mathrm{params}}
+ \mathrm{bucket}(M_{\mathrm{params}})
+ \min(M_{\mathrm{params}}, 256~\mathrm{MiB}),
\\
M_{\mathrm{zero,update}} &=
M_{\mathrm{opt}}
+ M_{\mathrm{grads}}
+ M_{\mathrm{update,scratch}},
\\
M_{\mathrm{zero,comm}} &=
M_{\mathrm{grads}}
+ \mathrm{bucket}(M_{\mathrm{grads}}).
\end{align}

\subsection{Phase-Local Workspace Terms}
The forward workspace is:
\begin{equation}
M_{\mathrm{ws,fwd}} =
M_{\mathrm{attn,fwd}}
+ M_{\mathrm{tp,comm}}
+ M_{\mathrm{sp,comm}}.
\end{equation}

The backward workspace is:
\begin{equation}
M_{\mathrm{ws,bwd}} =
M_{\mathrm{bwd,kernel}}
+ M_{\mathrm{recompute}}
+ M_{\mathrm{tp,comm}}
+ M_{\mathrm{sp,comm}}.
\end{equation}

The optimizer workspace depends on backend:
\begin{equation}
M_{\mathrm{ws,opt}} =
\begin{cases}
M_{\mathrm{update,scratch}} & \text{for single\_gpu}, \\
\max(M_{\mathrm{ddp,comm}}, M_{\mathrm{update,scratch}}) & \text{for ddp}, \\
\max(M_{\mathrm{zero,fetch}}, M_{\mathrm{zero,update}}, M_{\mathrm{zero,comm}})
& \text{for zero2 or zero3}.
\end{cases}
\end{equation}

\subsection{Backward End State}
The phase model carries a distinct backward-end state into optimizer step:
\begin{equation}
M_{\mathrm{bwd,end}} =
\begin{cases}
0 & \text{for single\_gpu}, \\
M_{\mathrm{grads}} + \mathrm{bucket}(M_{\mathrm{grads}}) & \text{for ddp}, \\
M_{\mathrm{grads}} + \mathrm{bucket}(M_{\mathrm{grads}}) & \text{for zero2}, \\
M_{\mathrm{grads}} + \mathrm{bucket}(M_{\mathrm{params}}) & \text{for zero3}.
\end{cases}
\end{equation}

\section{Phase Assembly}
The phase model constructs six synthetic records:
\begin{enumerate}[leftmargin=1.5em]
\item post-init baseline,
\item forward,
\item loss materialization,
\item backward,
\item optimizer step,
\item step end.
\end{enumerate}

The baseline is:
\begin{equation}
M_{\mathrm{baseline}} = M_{\mathrm{resident,base}}.
\end{equation}

The main peaks are:
\begin{align}
M_{\mathrm{forward}} &=
M_{\mathrm{resident,base}}
+ M_{\mathrm{act,fwd}}
+ M_{\mathrm{ws,fwd}},
\\
M_{\mathrm{backward}} &=
M_{\mathrm{resident,with\ grad}}
+ M_{\mathrm{act,bwd}}
+ M_{\mathrm{ws,bwd}},
\\
M_{\mathrm{optimizer}} &=
M_{\mathrm{resident,with\ grad}}
+ M_{\mathrm{bwd,end}}
+ M_{\mathrm{ws,opt}}.
\end{align}

The global peak phase is whichever of \texttt{forward}, \texttt{backward}, and
\texttt{optimizer\_step} has the largest allocated peak.

\section{Public Breakdown Semantics}
The public \texttt{MemoryComponentBreakdown} is aligned to the selected peak
phase.

\subsection{Always-Present Resident Terms}
\begin{align}
M_{\mathrm{breakdown,params}} &= M_{\mathrm{params}}, \\
M_{\mathrm{breakdown,grads}} &= M_{\mathrm{grads}}, \\
M_{\mathrm{breakdown,opt}} &= M_{\mathrm{opt}}, \\
M_{\mathrm{breakdown,runtime}} &= M_{\mathrm{runtime}} + M_{\mathrm{master}}.
\end{align}

\subsection{Phase-Dependent Activation and Transient Terms}
\paragraph{Forward peak.}
\begin{align}
M_{\mathrm{breakdown,act}} &= M_{\mathrm{act,fwd}}, \\
M_{\mathrm{breakdown,transient}} &= M_{\mathrm{ws,fwd}}.
\end{align}

\paragraph{Backward peak.}
\begin{align}
M_{\mathrm{breakdown,act}} &= M_{\mathrm{act,bwd}}, \\
M_{\mathrm{breakdown,transient}} &= M_{\mathrm{ws,bwd}}.
\end{align}

\paragraph{Optimizer-step peak.}
\begin{align}
M_{\mathrm{breakdown,act}} &= 0, \\
M_{\mathrm{breakdown,transient}} &= M_{\mathrm{bwd,end}} + M_{\mathrm{ws,opt}}.
\end{align}

This is an important semantic choice: optimizer step does not report retained
activation bytes. Any carried post-backward state is shown in the transient
term instead.

\section{Debug Metadata}
The estimate result carries typed debug structures:
\begin{itemize}[leftmargin=1.5em]
\item resident-state debug:
  parameter, gradient, optimizer-state, master-weight, runtime-floor, and
  trainable-parameter bytes,
\item activation debug:
  all explicit retained activation classes plus the hook-visible diagnostic,
\item workspace debug:
  all local workspace and communication windows,
\item phase-peak debug:
  forward, backward, optimizer, and global peaks.
\end{itemize}

\paragraph{Reserved-peak fields.}
The debug object still carries
\begin{equation}
M_{\mathrm{soft\_reserved}}, \qquad
M_{\mathrm{stressed\_reserved}},
\end{equation}
but in the current simplified implementation they are equal to the allocated
peak. They are placeholders for reporting compatibility, not an independent
allocator model in the main estimator path.

\section{What the Estimator Explicitly Does Not Model}
The current implementation is intentionally explicit about its limits.

\subsection{Node Topology}
Multi-node affects the estimator only through total available GPU count and the
derived DP degree. The estimator does \emph{not} distinguish:
\begin{itemize}[leftmargin=1.5em]
\item one node with eight GPUs,
\item two nodes with four GPUs,
\item eight nodes with two GPUs,
\end{itemize}
if the total world size and TP degree are the same. Inter-node bandwidth,
topology, and placement are not modeled in the memory equations.

\subsection{Runtime Measurement Coverage}
The estimator supports TP and SP analytically. The runtime measurement path does
not. Measurement currently rejects:
\begin{itemize}[leftmargin=1.5em]
\item \texttt{tensor\_parallel\_degree > 1},
\item \texttt{sequence\_parallel = True}.
\end{itemize}

\subsection{TP/SP Structural Simplifications}
The current TP/SP model is structural, not graph-exact:
\begin{itemize}[leftmargin=1.5em]
\item TP sharding is inferred from parameter categories and a small row-parallel
name heuristic,
\item SP is modeled as a boolean that reuses the TP group, not as an
independent degree,
\item TP and SP communication windows are simple hidden-state exchange proxies.
\end{itemize}

\subsection{Attention Backends}
Only \texttt{standard} attention receives an explicit retained score tensor.
\texttt{sdpa} and \texttt{flash2} are modeled without retained score storage.
Their memory pressure is expressed through the forward workspace and backward
kernel workspace instead.

\subsection{Optimizer Step Idealization}
The optimizer step is modeled as:
\begin{itemize}[leftmargin=1.5em]
\item a resolved per-parameter update scratch,
\item plus explicit DDP or ZeRO windows,
\item plus a backward-end state.
\end{itemize}

It does not trace exact kernel fusion, bucket overlap timing, or allocator
retention. The current ZeRO scratch sharding by world size is a deliberate
structural assumption, not a validated runtime trace.

\subsection{Loss Workspace}
The current \texttt{WorkspaceDebug} still has a
\texttt{loss\_workspace\_bytes} field, but the simplified estimator sets it to
zero.

\section{Bottom Line}
The current estimator can be summarized as:
\begin{equation}
\text{resident state}
+ \text{retained saved tensors}
+ \text{explicit local and communication windows}
\Longrightarrow
\text{synthetic phase peaks}.
\end{equation}

The important practical consequences are:
\begin{itemize}[leftmargin=1.5em]
\item retained activations are built from named tensor classes, not from
hook-visible outputs,
\item checkpointing moves rematerializable work into backward workspace,
\item TP and SP are modeled as orthogonal structural axes on top of
DDP/ZeRO-style distributed modes,
\item multi-node is represented through world size, but not through node-aware
communication topology,
\item optimizer-step peaks are explicit and inspectable, but still idealized.
\end{itemize}

That is the current implemented estimation approach.

\end{document}
